\chapter*{Cómo leer este libro}
\addcontentsline{toc}{chapter}{Cómo leer este libro}
\markboth{Cómo leer este libro}{Cómo leer este libro}

Cada capítulo abre con una lista de objetivos y cierra con un resumen, una serie de ejercicios
graduados y las referencias citadas. Entre medias, el texto se apoya en un pequeño repertorio de
recuadros con funciones precisas. Conocerlos de antemano le ayudará a decidir qué leer con calma
y qué puede saltarse en una primera lectura.

\begin{intuicion}[Recuadros de intuición]
Abren cada tema con una situación cotidiana que sirve de punto de apoyo. No sustituyen la teoría:
preparan el terreno para ella.
\end{intuicion}

\begin{definicion*}{Definiciones}{}
Enuncian con precisión los objetos matemáticos o biológicos con los que trabajaremos. Están
numeradas por capítulo.
\end{definicion*}

\begin{teorema*}{Resultados}{}
Recogen las ecuaciones y los resultados centrales de cada método, normalmente seguidos de una
tabla de símbolos.
\end{teorema*}

\begin{simbolos}
\simbolo{$x$}{Cada ecuación importante va acompañada de una tabla como esta, que explica el significado de todos sus símbolos.}
\end{simbolos}

\begin{ejemplo*}{Ejemplos resueltos}{}
Aplican la teoría a números concretos, paso a paso, para que usted pueda reproducirlos con lápiz
o con Python.
\end{ejemplo*}

\begin{codigo}[ejemplo.py]
# Los fragmentos de código son breves, completos y ejecutables.
gc = sum(base in "GC" for base in "ATGCGC") / 6
\end{codigo}

\begin{historia}[Notas históricas]
Cuentan quién, cuándo y por qué se desarrolló una idea. La ciencia es una actividad humana.
\end{historia}

\begin{cuidado}[Advertencias]
Señalan errores frecuentes, supuestos que se olvidan y trampas de interpretación.
\end{cuidado}

\begin{profundizacion}[Para profundizar]
Material avanzado que puede omitirse en una primera lectura sin perder el hilo.
\end{profundizacion}

\enlacerepo{enlace al laboratorio de cada lección}

\paragraph{Convenciones} Las explicaciones están en español; el código, los nombres de funciones
y muchos términos técnicos se mantienen en inglés, como aparecen en la literatura y en la
documentación de las herramientas; su primera aparición se marca en \ingles{cursiva}. Los
decimales se escriben con coma. Las secuencias de ADN se colorean con el código de la
\cref{tab:convenciones}, el mismo que se usa en los notebooks del curso.

\begin{table}[h]
\centering\small
\caption{Colores de nucleótidos usados en todo el libro.}\label{tab:convenciones}
\begin{tabular}{@{}cccc@{}}
\toprule
\nuc{A} adenina & \nuc{C} citosina & \nuc{G} guanina & \nuc{T} timina (\nuc{U} uracilo)\\
\bottomrule
\end{tabular}
\end{table}

\paragraph{Ejercicios} Están marcados según su dificultad:
\dificultad{1}~básica, \dificultad{2}~intermedia y \dificultad{3}~avanzada. Los de nivel avanzado
suelen requerir programación o una demostración.

\paragraph{Requisitos} Se asume una familiaridad básica con cálculo, probabilidad y programación en
Python. El \cref{cap:00} repasa las herramientas computacionales y el \cref{cap:01}, la biología
molecular necesaria.
